%%% SECTION 11: NATIONAL MCP SERVER INFRASTRUCTURE %%%

[This section should be approximately 10-12 pages. The major theme is that
Physical AI National MCP Servers \cite{national-mcp-paper} provides the
national-scale infrastructure for connecting Physical AI oncology trial
systems across sites using the Model Context Protocol \cite{mcp-protocol}.
The five-server architecture addresses the fragmentation problem that prevents
current clinical trial systems from interoperating, and establishes
standardized communication pathways for Physical AI operations at scale.]

[PRIMARY SOURCE: national-platform/national\_mcp/ directory containing four
chunked files: chunk1\_preamble\_and\_introduction.tex (lines 1-181),
chunk2\_methods\_and\_results\_part1.tex (lines 182-465),
chunk3\_results\_part2.tex (lines 466-689), and
chunk4\_discussion\_and\_conclusion.tex (lines 690-1011). All four files must
be used together to build this section.]

\subsection{The Fragmentation Problem}
\label{subsec:mcp-fragmentation}

[Build from national-platform/national\_mcp/chunk1\_preamble\_and\_introduction.tex,
the introduction section on fragmentation. Detail how current clinical trial
systems operate in silos with incompatible data formats, disconnected
communication channels, and no standardized protocol for real-time inter-system
coordination. Explain how this fragmentation is magnified in Physical AI
trials where robots, AI systems, clinical databases, and regulatory reporting
tools must all communicate seamlessly.]

[Reference the current trial infrastructure from ClinicalTrials.gov
\cite{clinicaltrials-gov} and FDA clinical trials guidance
\cite{fda-clinical-trials-guidance} to show the status quo. Explain how
FHIR \cite{hl7-fhir} and DICOM \cite{dicom} address part of the problem
but do not provide a complete solution for Physical AI systems.]

\subsection{Model Context Protocol for Physical AI}
\label{subsec:mcp-protocol}

[Build from national-platform/national\_mcp/chunk1\_preamble\_and\_introduction.tex
on the MCP opportunity and prior architecture. Explain how the Model Context
Protocol \cite{mcp-protocol} provides a standardized framework for AI
systems to communicate with tools and data sources, and how this framework is
adapted for Physical AI oncology trial operations. Detail the prior
architecture that the five-server design builds upon.]

\subsection{Five-Server Architecture}
\label{subsec:mcp-architecture}

[Build from national-platform/national\_mcp/chunk2\_methods\_and\_results\_part1.tex,
the five-server architecture section. Detail each of the five servers, their
roles, capabilities, and interrelationships. Include the tool inventories for
each server and explain how tools are organized to support different aspects of
Physical AI trial operations.]

[MAJOR THEME: The five-server architecture provides national-scale
infrastructure that can be replicated at each trial site while maintaining
centralized coordination. This architecture supports the goal of scaling
Physical AI oncology trials nationwide by providing standardized infrastructure
that any site can deploy.]

\subsection{Tool Inventories and Conformance Levels}
\label{subsec:mcp-tools}

[Build from national-platform/national\_mcp/chunk2\_methods\_and\_results\_part1.tex
on tool inventories and conformance levels. Detail the complete set of tools
available across all five servers, the conformance level requirements for each
tool category, and how conformance is verified. Reference the MCP repository
\cite{mcp-repo} for implementation details.]

\subsection{Safety Modules and Emergency Systems}
\label{subsec:mcp-safety}

[Build from national-platform/national\_mcp/chunk3\_results\_part2.tex on
safety modules. Cover the safety architecture embedded in the MCP server
infrastructure, including emergency stop capabilities, automatic failover,
safety interlock protocols, and the emergency stop code listing. Cross-reference
with the emergency preparedness requirements in
Section~\ref{sec:site-establishment}.]

\subsection{Integration and Interoperability}
\label{subsec:mcp-integration}

[Build from national-platform/national\_mcp/chunk3\_results\_part2.tex on
integration adapters, test coverage, deployment infrastructure, SDK/tooling,
and schemas. Cover the interoperability testbed and how it validates that
MCP servers can communicate with existing healthcare systems using FHIR
\cite{hl7-fhir} and DICOM \cite{dicom}.]

\subsection{National Deployment Topology}
\label{subsec:mcp-deployment}

[Build from national-platform/national\_mcp/chunk4\_discussion\_and\_conclusion.tex,
the national deployment topology section. Detail the planned deployment
architecture for nationwide Physical AI trial operations, including hub-and-spoke
models, regional server clusters, data routing, and redundancy provisions.]

\subsection{Governance and Compliance}
\label{subsec:mcp-governance}

[Build from national-platform/national\_mcp/chunk4\_discussion\_and\_conclusion.tex
on governance and comparison table. Cover the governance framework for national
MCP server operations, including who manages the servers, how updates are
deployed, compliance monitoring, and dispute resolution. Reference HIPAA
\cite{hipaa} for data governance requirements and the three adapted regulatory
standards for trial-specific governance.]

\subsection{CI/CD and Operational Excellence}
\label{subsec:mcp-cicd}

[Build from national-platform/national\_mcp/chunk4\_discussion\_and\_conclusion.tex
on CI/CD. Cover the continuous integration and deployment pipeline for MCP
servers, including automated testing, deployment procedures, rollback
capabilities, and version management. Reference the main repository
\cite{main-repo} and MCP repository \cite{mcp-repo} for implementation.]

\subsection{Patient Safety and Data Flow}
\label{subsec:mcp-patient-safety}

[Build from national-platform/national\_mcp/chunk4\_discussion\_and\_conclusion.tex
on patient safety and robot procedure data flow. Detail how MCP servers ensure
patient safety through real-time monitoring, data validation, and automated
safety checks. Cover the complete data flow from robot procedures through
MCP servers to clinical databases and regulatory reporting systems.]

[Connect to Physical AI Federated Learning \cite{fl-paper} and federated
learning in healthcare \cite{federated-learning-healthcare} for how MCP
servers interface with the federated learning framework to enable multi-site
coordination while preserving patient privacy.]
